\subsection{固有模空间的存在性}\label{src:9.3.2-subsection}
回顾 $\Theta$-完备性和 S-完备性是赋值判据 (定义~\bookref{7.9.10} 和 \bookref{7.9.14}), 它们保证分离良模空间的存在性 (定理~\bookref{7.10.1}). 我们将利用来自 $\Theta$ 的映射与凝聚层滤过之间的解释 (命题~\bookref{7.9.6}), 以及参数化凝聚层的叠 $\Coh(X)$ 是 $\Theta$-完备且 S-完备的事实 (命题~\bookref{7.9.30}), 验证 $\cBun_{r,d}^{\mathrm{ss}}$ 是 $\Theta$-完备且 S-完备的. 后一事实对任意射影概形 $X$ 成立. 为此, 给出 $\Coh(X)$ 的开子叠是 $\Theta$-完备或 S-完备的一般判据. 为表述 S-完备性判据, 需要如下概念: 两个 $\mathbb Z$-分次滤过
\[
E_\bullet:0\subseteq\cdots\subseteq E_{i-1}\subseteq E_i\subseteq E_{i+1}\subseteq\cdots\subseteq E
\]
和
\[
F^\bullet:F\supseteq\cdots\supseteq F^{i-1}\supseteq F^i\supseteq F^{i+1}\supseteq\cdots\supseteq0
\]
称为\newterm{相反的}{opposite}, 若对所有 $i$, 有 $E_i/E_{i-1}\cong F^i/F^{i+1}$. 推广滤过与来自 $\Theta$ 的映射之间的对应, 满足 $i\ll0$ 时 $E_i=0$, $F^i=F$, 且 $i\gg0$ 时 $E_i=E$, $F^i=0$ 的相反滤过 $E_\bullet,F^\bullet$ 的数据, 等价于给出映射
\[
[\Spec(\Bbbk[x,y]/xy)/\mathbb G_m]\to\Coh(X),
\]
其中 $t\cdot(x,y)=(tx,t^{-1}y)$. 此时该映射将 $(1,0)$ 映到 $E$, 将 $(0,1)$ 映到 $F$, 将 $(0,0)$ 映到 $\operatorname{gr}E_\bullet\cong\operatorname{gr}F^\bullet$.
\begin{proposition}\label{src:9.3.6}
设 $X$ 是代数闭域 $\Bbbk$ 上的射影概形, $\mathcal U\subseteq\Coh(X)$ 是开子叠.
\begin{enumerate}[(1)]
\item 子叠 $\mathcal U$ 是 $\Theta$-完备的, 当且仅当对每个 DVR $R$ (分式域为 $K$, 剩余域为 $\kappa$), 每个 $X_R$ 上在 $R$ 上平坦的凝聚层 $\mathcal E$, 以及每个满足 $i\ll0$ 时 $\mathcal E_i=0$, $i\gg0$ 时 $\mathcal E_i=\mathcal E$, 且每个 $\mathcal E_i/\mathcal E_{i-1}$ 在 $R$ 上平坦的 $\mathbb Z$-分次滤过 $\mathcal E_\bullet$, 若 $\mathcal E$ 和 $\operatorname{gr}(\mathcal E_\bullet|_K)$ 在 $\mathcal U$ 中, 则 $\operatorname{gr}(\mathcal E_\bullet|_\kappa)$ 也必须在 $\mathcal U$ 中.
\item 若对每个域扩张 $\kappa/\Bbbk$ 以及 $E,F\in\mathcal U(\kappa)$ 的每对相反滤过 $E_\bullet,F^\bullet$, 伴随分次 $\operatorname{gr}E_\bullet$ 在 $\mathcal U(\kappa)$ 中, 则子叠 S-完备.
\end{enumerate}
\end{proposition}
\begin{proof}
因为已知 $\Coh(X)$ 是 $\Theta$-完备且 S-完备的 (命题~\bookref{7.9.30}), $\mathcal U$ 的赋值判据等价于对所有如下交换图存在提升:
\[
\begin{tikzcd}[ampersand replacement=\&,column sep=small]
\Theta_R\setminus0\arrow[r]\arrow[d,hook]\&\mathcal U\arrow[d,hook]\&\Phi_R\setminus0\arrow[r]\arrow[d,hook]\&\mathcal U\arrow[d,hook]\\
\Theta_R\arrow[r]\arrow[ur,dashed]\&\Coh(X)\&\Phi_R\arrow[r]\arrow[ur,dashed]\&\Coh(X),
\end{tikzcd}
\]
其中 $R$ 是 DVR. 换言之, 需说明在唯一的填充 $\Theta_R\to\Coh(X)$ 和 $\Phi_R\to\Coh(X)$ 下, $0$ 的像包含于 $\mathcal U$. 由命题~\bookref{7.9.6}, 映射 $\Theta_R\to\Coh(X)$ 由 $C_R$ 上在 $R$ 上平坦的凝聚层 $\mathcal E$, 及满足 $i\ll0$ 时 $\mathcal E_i=0$, $i\gg0$ 时 $\mathcal E_i=\mathcal E$, 且每个 $\mathcal E_i/\mathcal E_{i-1}$ 在 $R$ 上平坦的 $\mathbb Z$-分次滤过 $\mathcal E_\bullet$ 分类. 此外, 限制 $\Theta_R\setminus0\to\Coh(X)$ 经过 $\mathcal U$ 分解, 当且仅当 $\mathcal E$ 和 $\operatorname{gr}(\mathcal E_\bullet|_K)$ 在 $\mathcal U$ 中. 因为 $0$ 在 $\Theta_R\to\Coh(X)$ 下的像为 $\operatorname{gr}(\mathcal E_\bullet|_\kappa)$, 得到 (1).

对 (2), 将 $\Phi_R\to\Coh(X)$ 沿 $\pi=0$ 限制, 得到对应于相反滤过 $E_\bullet,F^\bullet$ 的映射 $[\Spec(\kappa[x,y]/xy)/\mathbb G_m]\to\Coh(X)$. 若 $\operatorname{gr}E_\bullet$ 在 $\mathcal U$ 中, 则 $\Phi_R\to\Coh(X)$ 的像包含于 $\mathcal U$. 另见 \cite[命题 3.17]{WZ25}.
\end{proof}
虽然后文不需要, 下述习题在概念上有帮助.
\begin{exercise}\label{src:9.3.7}
证明当 $X=C$ 是曲线时, 命题~\bookref{9.3.6}(2) 的逆命题成立.

\textit{提示:} 令 $\mathcal X_0:=[\Spec(\Bbbk[x,y]/xy)/\mathbb G_m]\subset\Phi_R$, $\mathcal X_n$ 为其幂零加厚. 利用形变理论, 将映射 $\mathcal X_0\to\Coh(C)$ 延拓为相容映射 $\mathcal X_n\to\Coh(C)$, 再利用凝聚 Tannaka 对偶 (\bookref{7.6.9}) 进一步延拓为映射 $\Phi_R\to\Coh(C)$.
\end{exercise}
\begin{proposition}\label{src:9.3.8}
对代数闭域 $\Bbbk$ 上的光滑连通射影曲线 $C$, 代数叠 $\cBun_{r,d}^{\mathrm{ss}}$ 是 $\Theta$-完备且 S-完备的.
\end{proposition}
\begin{proof}
应用命题~\bookref{9.3.6} 建立这两种性质. 对 $\Theta$-完备性, 设 $\mathcal E$ 是 $C_R$ 上在 $R$ 上平坦的凝聚层, $\mathcal E_\bullet$ 是 $\mathbb Z$-分次滤过, 满足 $i\ll0$ 时 $\mathcal E_i=0$, $i\gg0$ 时 $\mathcal E_i=\mathcal E$, 且每个 $\mathcal E_i/\mathcal E_{i-1}$ 在 $R$ 上平坦. 假设 $\mathcal E$ 和 $\operatorname{gr}(\mathcal E_\bullet|_K)$ 在 $\cBun_{r,d}^{\mathrm{ss}}$ 中. 因为次数与秩在平坦族中不变, $\operatorname{gr}(\mathcal E_\bullet|_\kappa)\in\Coh_{r,d}$. 为说明 $\operatorname{gr}(\mathcal E_\bullet|_\kappa)$ 半稳定, 只需说明每个因子 $(\mathcal E_i/\mathcal E_{i-1})_\kappa$ 半稳定且斜率为 $d/r$. 因为 $(\mathcal E_i)_\kappa\subseteq\mathcal E_\kappa$ 是半稳定层中斜率相同的子层, 每个 $(\mathcal E_i)_\kappa$ 半稳定. 由于固定斜率的半稳定层组成 Abel 范畴 (习题~\bookref{9.1.17}), $(\mathcal E_i/\mathcal E_{i-1})_\kappa=(\mathcal E_i)_\kappa/(\mathcal E_{i-1})_\kappa$ 半稳定且斜率为 $d/r$.

对 S-完备性, 设
\[
\begin{aligned}
E_\bullet&:\cdots\subseteq E_{i-1}\subseteq E_i\subseteq\cdots\subseteq E,\\
F^\bullet&:F\supseteq\cdots\supseteq F^i\supseteq F^{i+1}\supseteq\cdots
\end{aligned}
\]
是域扩张 $\kappa/\Bbbk$ 上 $C_\kappa$ 上半稳定向量丛 $E,F$ 的相反滤过. 这意味着 $E_i/E_{i-1}\cong F^i/F^{i+1}$. 归纳证明每个 $E_i,F^i$ 半稳定且斜率为 $d/r$. 由习题~\bookref{9.1.17}, 这建立所需的 $\operatorname{gr}E_\bullet=\bigoplus E_i/E_{i-1}$ 的半稳定性. 因为 $i\ll0$ 时 $E_i=0$, $F_i=F$, 且 $F$ 半稳定, $i\ll0$ 时 $E_j,F_j$ 半稳定. 假设 $j\leq i$ 时 $E_j,F_j$ 半稳定. 因为 $F_{i+1}\cong\ker(F_i\to E_i/E_{i-1})$ 是同斜率半稳定向量丛之间映射的核, $F_{i+1}$ 半稳定且斜率为 $d/r$. 另一方面, $E_{i+1}/E_i\cong F_{i+1}/F_{i+2}$ 既是半稳定向量丛 $E/E_i$ 的子层 (利用 $E$ 的半稳定性和归纳假设), 又是半稳定向量丛 $F_{i+1}$ 的商. 所以 $\mu(E_{i+1}/E_i)=d/r$, 从而 $E_{i+1}$ 半稳定且斜率为 $d/r$.
\end{proof}
向量丛叠 $\cBun_{r,d}$ 则既不 $\Theta$-完备, 也不 S-完备.
\begin{example}\label{src:9.3.9}
设 $g\geq2$. 令 $p\in C(\Bbbk)$ 为截面 $s\in\Gamma(C,\mathcal O(p))$ 的零点, $I\subseteq\mathcal O_{C_R}$ 为 $(p,0)\in C\times\Spec R$ 的理想层. 单射 $(s,-\pi):\mathcal O_{C_R}(-p)\hookrightarrow\mathcal O_{C_R}\oplus\mathcal O_{C_R}(-p)$ 的商为 $I$, 它无挠, 因而在 $R$ 上平坦, 但不是向量丛. 由命题~\bookref{9.3.6}, $\cBun_{2,-1}$ 不 $\Theta$-完备.

设 $L,M$ 为 $C$ 上线丛, $p\in C(\Bbbk)$ 是使 $\operatorname{Ext}_{\mathcal O_C}^1(M,L(p))$ 和 $\operatorname{Ext}_{\mathcal O_C}^1(L,M(p))$ 非零的点; 若 $L,M$ 次数相同, Riemann--Roch 计算说明这两个 $\operatorname{Ext}^1$ 群都非零. 令 $Q$ (分别地, $Q'$) 是 $M$ 被 $L(p)$ (分别地, $L$ 被 $M(p)$) 的非平凡扩张. 则
\[
E_\bullet:0\subseteq L\subseteq L(p)\subseteq Q
\quad\text{和}\quad F^\bullet:Q'\supseteq M(p)\supseteq M\supseteq0
\]
定义相反滤过, 其中 $E_0=L$, $F^0=Q'$. 伴随分次 $\operatorname{gr}E_\bullet=L\oplus\kappa(p)\oplus M$ 不是向量丛. 由命题~\bookref{9.3.6}, $\cBun_{1,d}$ 不 S-完备, 其中 $d=\deg L+\deg M+1$.
\end{example}
\begin{exercise}\label{src:9.3.10}
证明 $\cBun_{r,d}^{\mathrm{simple}}$ 不 S-完备, 并推出 $M_{r,d}^{\mathrm{simple}}$ 不分离. (由习题~\bookref{9.2.25}, $\cBun_{r,d}^{\mathrm{simple}}$ 和 $M_{r,d}^{\mathrm{simple}}$ 在 $\Bbbk$ 上都不是有限型的.)
\end{exercise}
\begin{theorem}\label{src:9.3.11}
设 $C$ 是特征为 $0$ 的代数闭域 $\Bbbk$ 上的光滑连通射影曲线. 对每个 $r>0$ 和 $d$, 存在 Cartesian 图
\[
\begin{tikzcd}[ampersand replacement=\&]
\cBun_{r,d}^{\mathrm s}\arrow[r]\arrow[d]\&\cBun_{r,d}^{\mathrm{ss}}\arrow[d]\\
M_{r,d}^{\mathrm s}\arrow[r]\&M_{r,d}^{\mathrm{ps}},
\end{tikzcd}
\]
其中 $M_{r,d}^{\mathrm{ps}}$ 是 $\Bbbk$ 上的固有代数空间, $\cBun_{r,d}^{\mathrm{ss}}\to M_{r,d}^{\mathrm{ps}}$ 是良模空间, $\cBun_{r,d}^{\mathrm s}\to M_{r,d}^{\mathrm s}$ 是带 $\mathbb G_m$-带的 gerbe. 此外,
\[
\dim M_{r,d}^{\mathrm{ps}}=\dim\cBun_{r,d}^{\mathrm{ss}}+1=r^2(g-1)+1.
\]
\end{theorem}
\begin{proof}
由半稳定性的有界性 (\bookref{9.2.8}), $M_{r,d}^{\mathrm{ps}}$ 有限型, 且它是 $\Theta$-完备且 S-完备的 (命题~\bookref{9.3.8}); 因而分离代数空间 $\pi:\cBun_{r,d}^{\mathrm{ss}}\to M_{r,d}^{\mathrm{ps}}$ 的存在性直接来自良模空间存在定理 (\bookref{7.10.1}). 由于 $\cBun_{r,d}^{\mathrm{ss}}$ 由 Langton 定理 (\bookref{9.3.5}) 普遍闭, $M_{r,d}^{\mathrm{ps}}$ 固有.

因为每个稳定丛多重稳定, 从而由命题~\bookref{9.1.35} 为 $\cBun_{r,d}^{\mathrm{ss}}$ 的闭点, 得到 $\pi^{-1}(\pi(\cBun_{r,d}^{\mathrm s}))=\cBun_{r,d}^{\mathrm s}$. 由于 $\pi$ 普遍闭, $\pi(\cBun_{r,d}^{\mathrm s})\subseteq M_{r,d}^{\mathrm{ps}}$ 开. 另一方面, 推论~\bookref{9.2.26} 给出带 $\mathbb G_m$-带的 gerbe $\cBun_{r,d}^{\mathrm s}\to M_{r,d}^{\mathrm s}$ 的存在性. 良模空间与带 $\mathbb G_m$-带的 gerbe 都对到代数空间的映射万有, 因而 $M_{r,d}^{\mathrm s}\cong\pi(\cBun_{r,d}^{\mathrm s})$.

为计算维数, 使用 $M_{r,d}^{\mathrm s}\subseteq M_{r,d}^{\mathrm{ps}}$ 开且稠密; 这来自 $\cBun_{r,d}^{\mathrm s}\subseteq\cBun_{r,d}$ 开且稠密的事实 (定理~\bookref{9.2.16}). 因为 $\cBun^{\mathrm s}\to M_{r,d}^{\mathrm s}$ 是带 $\mathbb G_m$-带的 gerbe, 且 $\dim B\mathbb G_m=-1$, 得到
\[
\dim M_{r,d}^{\mathrm{ps}}=\dim M_{r,d}^{\mathrm s}
=\dim\cBun_{r,d}^{\mathrm s}+1=r^2(g-1)+1,
\]
其中利用了命题~\bookref{4.7.9} 对 $\cBun_{r,d}$ 维数的计算.
\end{proof}
\begin{exercise}\label{src:9.3.12}
对每个 $\xi\in\mathrm{Pic}(C)$, 以 $\cBun_{r,\xi}^{\mathrm{ss}}$ (分别地, $\cBun_{r,\xi}^{\mathrm s}$) 表示模叠 $\cBun_{r,\xi}$ 中参数化行列式为 $\xi$ 的半稳定 (分别地, 稳定) 向量丛的子叠 (定义~\bookref{9.2.27}). 证明 $\cBun_{r,\xi}^{\mathrm{ss}}$ 是 $\Theta$-完备且 S-完备的. 推出在特征 $0$ 下存在固有良模空间 $\cBun_{r,\xi}^{\mathrm{ss}}\to M_{r,\xi}^{\mathrm{ps}}$, 限制为 $\mu_r$-gerbe $\cBun_{r,\xi}^{\mathrm s}\to M_{r,\xi}^{\mathrm s}$, 且 $M_{r,\xi}^{\mathrm s}\subseteq M_{r,\xi}^{\mathrm{ps}}$ 为开子空间.

\textit{提示:} 利用 $\cBun_{r,\xi}^{\mathrm{ss}}\to\cBun_{r,\deg\xi}^{\mathrm{ss}}$ 仿射.
\end{exercise}
\begin{exercise}\label{src:9.3.13}
证明存在交换图
\[
\begin{tikzcd}[ampersand replacement=\&]
\cBun_{r,d}^{\mathrm{ss}}\arrow[r,"\det"]\arrow[d]\&\underline{\mathcal P\!ic}^d\arrow[d]\\
M_{r,d}^{\mathrm{ps}}\arrow[r,"\det"']\&\underline{\mathrm{Pic}}^d,
\end{tikzcd}
\]
其中 $\underline{\mathcal P\!ic}^d$ 和 $\underline{\mathrm{Pic}}^d$ 为 Picard 叠与 Picard 概形 (见 \S\bookref{7.4.8-subsection}). 对次数为 $d$ 的线丛 $\xi$, 证明 $M_{r,d}^{\mathrm{ps}}\to\underline{\mathrm{Pic}}^d$ 在 $\xi$ 上的纤维与 $M_{r,\xi}^{\mathrm{ps}}$ 等同.
\end{exercise}
注意 $\cBun_{r,d}^{\mathrm{ss}}(C)$ (同样也包括 $\cBun_{r,\xi}^{\mathrm{ss}}(C)$) 的每个闭点都有约化自同构群; 这可由 $\cBun_{r,d}^{\mathrm{ss}}(C)$ 的 S-完备性和命题~\bookref{7.10.30} 得到, 也可在特征 $0$ 下由良模空间的存在性和习题~\bookref{7.5.27} 得到. 因为 $\cBun_{r,d}^{\mathrm{ss}}(C)$ 的闭点对应于多重稳定向量丛 (这一事实一般成立, 但命题~\bookref{9.1.35} 仅在特征 $0$ 下证明), 再次得到多重稳定向量丛具有约化自同构群这一事实, 我们在推论~\bookref{9.1.16} 中已直接建立它. 严格半稳定但非多重稳定的向量丛则可能具有非约化自同构群.
\begin{exercise}\label{src:9.3.14}
构造具有非约化自同构群的半稳定向量丛.

\textit{提示:} 考虑稳定向量丛的非分裂扩张.
\end{exercise}
